\chapter{介词}

介词分组学，按格分，按话题分。

\section{按支配格分}
\begin{center}
    \begin{tabular} {|c|c|c|c|}
        \hline
        \olive[border=true] & \diam[border=true] & \olive[border=true] \& \diam[border=true] & \tri[border=true] \\
        \hline
        \makecell{mit                                                                                            \\ zu \\ aus \\ von \\ nach \\ bei \\ seit} & \makecell{durch \\ für \\ gegen \\ ohne \\ um \\ per} & \makecell{an \\ in \\ vor \\ ab \\ auf} & \makecell{inklusive \\ aufgrund} \\
        \hline
    \end{tabular}
\end{center}

\section{去}

\begin{itemize}
    \item auf: 去机关 \diam[border=true] (去用宾格，在用与格)
          \begin{itemize}
              \item auf die Party gehen
              \item auf die Post gehen
          \end{itemize}
    \item zu: 去小地方 \olive[border=true]
          \begin{itemize}
              \item zur Party kommen
              \item zum Bahnhof fahren
          \end{itemize}
    \item durch: 穿过 \diam[border=true]
          \begin{itemize}
              \item durch die Zollkontrolle gehen
          \end{itemize}
\end{itemize}

\section{in}
in \diam[border=true] gehen
\begin{itemize}
    \item ins Kino gehen
    \item ins Theater gehen
    \item ins Konzert gehen
    \item ins Bett gehen
    \item in den Park gehen
    \item in die Oper gehen
    \item in die Bibliothek gehen
    \item in die Schweiz gehen (仅去非中性国家用in)
\end{itemize}

in \olive[border=true] sein/bleiben
\begin{itemize}
    \item im Bett bleiben
    \item im Büro sein
    \item im Supermarkt kaufen
\end{itemize}

in \olive[border=true] 长时段\textsuperscript{\hyperref[time:prep:in]{见时间词汇}}

\section{an}
an \olive[border=true] 短时段\textsuperscript{\hyperref[time:prep:an]{见时间词汇}} \\
an 地点边上 sein
\begin{itemize}
    \item am Neckar sein
\end{itemize}

\newgeometry{left=3cm,right=8cm,marginparwidth=6cm}
\setlength{\parindent}{0pt}
\section{方位介词}


\subsection{来}

\subsubsection{aus \olive[border=true]}
\begin{itemize}
    \item 从某地里来，封闭 \\
          Er trinkt Bier aus einer Flasche. \marginpar{Flasche 封闭容器} \\
          Schau mal aus dem Fenster. \\
          Er kommt jetzt aus dem Klassenzimmer. \\
          Ich komme aus China. \\
    \item 由材质制成，(材质抽象所以零冠) \\
          Die Tasche ist aus Leder.
\end{itemize}

\subsubsection{von \olive[border=true]}
\begin{itemize}
    \item 从某方来，不封闭，可接抽象概念 \\
          Sie isst Salad vom Teller. \marginpar{Teller 不封闭容器} \\
          Mein Vater kommt sehr spät von der Arbeit. \\
          Er kommt vom seinem Freund. \\
\end{itemize}

\subsubsection{辨析：aus vs von}
\begin{itemize}
    \item 类比英文，aus $\approx$ out of, von $\approx$ from
    \item 可视化图解，aus是从圆里出来，von是直箭头
    \item 根据句义、逻辑、常识推断封闭性 \\
          朋友那，而不是朋友家
    \item 根据搭配 \\
          Sie ist gestern von Frankfurt nach Shanghai geflogen. \\
          von .. nach .. 强调方向而不是出来
\end{itemize}


\subsection{去}

\subsubsection{nach \olive[border=true]}
\begin{itemize}
    \item 去某大地方 (国、城)
    \item 去方位，东南西北，上下左右 \\
          Gehen Sie an der nächsten Kreuzung nach rechts. \\

\end{itemize}

\subsubsection{zu \olive[border=true]}
\begin{itemize}
    \item 去某小地方 (店、家)
    \item 去方向，接抽象名词 \\
          Er geht jeden Tag sehr Früh zur Arbeit.
\end{itemize}

\subsubsection{辨析：nach vs zu}
\begin{itemize}
    \item 类比英文，nach $\approx$ after, zu $\approx$ to
\end{itemize}

\subsection{Haus}
\begin{itemize}
    \item at home: zu Haus \\
          Er will den ganzen Tag zu Haus bleiben. \marginpar{at home 做状语} \\
          Vor einer Stunde ist Hans von zu Haus in Büro gekommen. \marginpar{von zu Haus里zu Haus是抽象概念} \\
    \item (go) home: nach Haus (gehen)
\end{itemize}

\subsection{在}
bei $\approx$ by
\subsubsection{bei \olive[border=true]}
\begin{itemize}
    \item 在某处 \\
          Sie sind beim Arzt. \\
          Monika arbeitet bei Siemens. \\
          Der Student wohnt bei seinen Eltern. \marginpar{bei jmdm wohnen: 住某人处 \\ 辨析 mit jmdm wohnen: 和某人住}
\end{itemize}
\subsubsection{gegenüber \olive[border=true]}
gegenüber $\approx$ across from
\begin{itemize}
    \item 在对面，常后置 \\
          Er wohnt mir gegenüber. (gegenüber 宾若为人称代词，则必后置) \\
          Das Kaufhaus liegt der Post gegenüber. \\
          Das Kaufhaus liegt gegenüber der Post. \\
\end{itemize}

\subsection{宾格}
\subsubsection{durch \diam[border=true]}
durch $\approx$ through
\begin{itemize}
    \item 穿过 \\
          Er geht durch den Wald.  \\
          Er reise durch Deutschland. \marginpar{travel through} \\
\end{itemize}

\subsubsection{gegen \diam[border=true]}
gegen $\approx$ toward
\begin{itemize}
    \item 朝着 \\
          Er läuft gegen die Wand.
\end{itemize}

\subsubsection{um \diam[border=true]}
um $\approx$ around
\begin{itemize}
    \item 绕着 \\
          Das Auto fährt um die Ecke.
    \item 围着 \\
          Wir sitzen um den Tisch.
\end{itemize}

\subsubsection{gegen, um 时间义项}
\begin{itemize}
    \item gegen: about 表大约几点
    \item um: at 表在几点
\end{itemize}

\subsection{静三动四 wohin? oder wo?}
\subsubsection{九介词}
静三动四九介词，都表方位。
\begin{center}
    \begin{tabular}{|c|c|c|c|}
        \hline
        \rowcolor{gray}
        direction & prep.          & eng equiv & adv.         \\
        \hline
        x         & \makecell{an                              \\ neben \\ zwischen}       & \makecell{at \\ next to \\ between} & \makecell{links \\ rechts}           \\
        \hline
        y         & \makecell{vor                             \\ hinter}  & \makecell{before \\ behind}        & \makecell{vorne \\ hinten}           \\
        \hline
        z         & \makecell{über                            \\ auf \\ unter} & \makecell{over \\ on \\ under} & \makecell{oben \\ unten} \\
        \hline
        o         & in             & in        & in der Mitte \\
        \hline
    \end{tabular}
\end{center}

注意：其他非方位义项可能只支配一格
\begin{itemize}
    \item warten auf \diam[border=true]
\end{itemize}

\subsubsection{方位动词}
\begin{center}
    \begin{tabular} {|c|c|c|c|c|}
        \hline
        \rowcolor{gray}
        word                      & meaning        & present & past         & perf      \\
        \hline
        setzen                    & v/t.: to sit   &         &              &           \\
        sizten                    & v/i.: to sit   &         & saß          & gesessen  \\
        \hline
        legen                     & v/t.: to lay   &         &              &           \\
        liegen                    & v/i.: to lie   &         & lag          & gelegen   \\
        \hline
        stellen                   & v/t.: to put   &         &              &           \\
        stehen                    & v/i.: to stand &         & stand        & gestanden \\
        \hline
        hängen                    & v/t.: to hang  &         & stand        & gestanden \\
        hängen\textsuperscript{2} & v/i.: to hang  &         & hing         & gehangen  \\
        \hline
        stecken                   & v/t.: to stuck &         &              &           \\
        stecken                   & v/i.: to stuck &         & steckte/stak &           \\
        \hline
    \end{tabular}
\end{center}

方位动词搭配介词
\begin{itemize}
    \item 及物为动，不及物为静
    \item 及物不规则，不及物规则
\end{itemize}


\restoregeometry

\section{介冠合写}
